% Part: normal-modal-logic
% Chapter: completeness
% Section: lindenbaums-lemma

\documentclass[../../../include/open-logic-section]{subfiles}

\begin{document}

\olfileid{nml}{com}{lin}

\olsection{લિન્ડનબાઉમનું સહાયક પ્રમેય}

લિન્ડનબાઉમનું સહાયક પ્રમેય સ્થાપિત કરે છે કે !!{formula}sનો દરેક
$\Sigma$-સુસંગત ગણ ઓછામાં ઓછા એક \emph{પૂર્ણ} $\Sigma$-સુસંગત ગણમાં
સમાયેલો છે. કૅનોનિકલ નિદર્શનની આપણી રચના બતાવશે કે દરેક પૂર્ણ
$\Sigma$-સુસંગત ગણ~$\Delta$ માટે કૅનોનિકલ નિદર્શનમાં એવું જગત છે
જ્યાં $\Delta$માંનાં બધાં અને માત્ર એ જ !!{formula}s સત્ય છે. તેથી
લિન્ડનબાઉમનું સહાયક પ્રમેય ખાતરી આપે છે કે દરેક $\Sigma$-સુસંગત ગણ
કૅનોનિકલ નિદર્શનના કોઈ જગત પર સત્ય છે.

\begin{thm}[લિન્ડનબાઉમનું સહાયક પ્રમેય]\ollabel{thm:lindenbaum}
  જો $\Gamma$ એ $\Sigma$-સુસંગત હોય, તો $\Gamma$ને વિસ્તારે એવો
  પૂર્ણ $\Sigma$-સુસંગત ગણ $\Delta$ છે.
\end{thm}

\begin{proof}
ભાષાનાં બધાં સૂત્રોની સંપૂર્ણ યાદી $!A_0$, $!A_1$, \dots{} લો
(પુનરાવર્તન માન્ય છે). દાખલા તરીકે, પહેલાં $\Obj p_0$ને યાદીમાં
મૂકો અને દરેક તબક્કે $n \ge 1$ માટે માત્ર $\Obj p_0$,
\dots,~$\Obj p_n$માંથી ચલ વાપરીને બનતાં લંબાઈ~$n$નાં સાન્ત
સંખ્યામાં સૂત્રોને યાદીમાં મૂકો. આપણે~$n$ પર અનુમાનથી !!{formula}sના
ગણો $\Delta_n$ વ્યાખ્યાયિત કરીએ છીએ અને પછી $\Delta = \bigcup_n
\Delta_n$ મૂકીએ છીએ. પહેલાં $\Delta_0 = \Gamma$ મૂકો. $\Delta_n$
વ્યાખ્યાયિત છે એમ ધારીને $\Delta_{n+1}$ આ રીતે વ્યાખ્યાયિત કરો:
\[
\Delta_{n+1} =
\begin{cases}
  \Delta_n \cup \{!A_n\}, & \text{જો $\Delta_n \cup \{ !A_n\}$
    એ $\Sigma$-સુસંગત હોય;} \\
  \Delta_n \cup \{ \lnot !A_n\}, & \text{અન્યથા.}
\end{cases}
\]
હવે $\Delta = \bigcup_{n=0}^\infty \Delta_n$ લો.

આ વ્યાખ્યાથી ખરેખર જરૂરી ગુણધર્મો ધરાવતો ગણ $\Delta$ મળે છે એવું
બતાવવું પડશે, એટલે કે $\Gamma \subseteq \Delta$ અને $\Delta$ એ
પૂર્ણ $\Sigma$-સુસંગત છે.

$\Gamma \subseteq \Delta$ સ્પષ્ટ છે, કારણ કે રચના મુજબ $\Delta_0
\subseteq \Delta$ અને $\Delta_0 = \Gamma$. હકીકતમાં, દરેક~$n$ માટે
$\Delta_n \subseteq \Delta$, કેમ કે $\Delta$ બધાં~$\Delta_n$નો સંઘ
છે. (રચનાના દરેક પગલે આપણે અગાઉના ગણમાં !!a{formula} ઉમેરીએ છીએ,
તેથી $\Delta_n \subseteq \Delta_{n+1}$; અને $\subseteq$ પરંપરિત
હોવાથી $n \le m$ હોય ત્યારે $\Delta_n \subseteq \Delta_{m}$.)
રચનાના દરેક તબક્કે આપણે કાં તો $!A_n$ અથવા $\lnot !A_n$ ઉમેરીએ
છીએ, અને દરેક !!{formula} બધાં~$!A_n$ની યાદીમાં (ઓછામાં ઓછું એક વાર)
આવે છે. તેથી દરેક $!A$ માટે કાં તો $!A \in \Delta$ અથવા
$\lnot !A \in \Delta$; એટલે વ્યાખ્યા મુજબ $\Delta$ પૂર્ણ છે.

છેલ્લે, $\Delta$ એ $\Sigma$-સુસંગત છે તે બતાવવું પડશે. આ માટે
બતાવીએ છીએ કે (a) જો $\Delta$ એ $\Sigma$-અસુસંગત હોય, તો કોઈ
$\Delta_n$ એ $\Sigma$-અસુસંગત હશે; અને (b) બધાં $\Delta_n$ એ
$\Sigma$-સુસંગત છે.

ધારો કે $\Delta$ એ $\Sigma$-અસુસંગત છે. ત્યારે $\Delta
\Proves[\Sigma] \lfalse$, એટલે કે એવાં $!A_1$, \dots,~$!A_k \in
\Delta$ છે કે $\Sigma \Proves !A_1 \lif (!A_2 \lif \cdots (!A_k
\lif \lfalse)\dots)$. $\Delta = \bigcup_{n=0}^\infty \Delta_n$
હોવાથી દરેક $!A_i \in \Delta_{n_i}$ કોઈક~$n_i$ માટે સાચું છે.
આ સૂચકોમાં સૌથી મોટો $n$ લો. $n_i \le n$ હોવાથી $\Delta_{n_i}
\subseteq \Delta_n$. આમ, બધાં $!A_i$ કોઈ એક~$\Delta_n$માં છે.
આનો અર્થ $\Delta_n \Proves[\Sigma] \lfalse$, એટલે કે $\Delta_n$
એ $\Sigma$-અસુસંગત છે.

દરેક $\Delta_n$ એ $\Sigma$-સુસંગત છે તે બતાવવા~$n$ પર સરળ અનુમાન
વાપરીએ. $\Delta_0 = \Gamma$, અને આપણે $\Gamma$ એ
$\Sigma$-સુસંગત છે એમ ધાર્યું હતું. એટલે $n = 0$ માટે દાવો સાચો
છે. હવે $n$ માટે દાવો સાચો છે, એટલે કે $\Delta_n$ એ
$\Sigma$-સુસંગત છે, એમ ધારો. જો $\Delta_n \cup \{!A_n\}$ એ
$\Sigma$-સુસંગત હોય, તો $\Delta_{n+1}$ એ તે ગણ છે; નહિ તો તે
$\Delta_n \cup \{\lnot!A_n\}$ છે. પહેલા કિસ્સામાં
$\Delta_{n+1}$ એ સ્પષ્ટ રીતે $\Sigma$-સુસંગત છે. પરંતુ
\olref[prf][con]{prop:consistencyfacts}\olref[prf][con]{prop:consistencyfacts-c}
મુજબ $\Delta_n \cup \{!A_n\}$ અથવા $\Delta_n \cup \{\lnot!A_n\}$
બેમાંથી ઓછામાં ઓછો એક સુસંગત છે; તેથી બીજા કિસ્સામાં પણ
$\Delta_{n+1}$ સુસંગત છે.
\end{proof}

\begin{cor}\ollabel{cor:provability-characterization}
  $\Gamma \Proves[\Sigma] !A$ જો અને માત્ર જો $\Gamma$ને વિસ્તારે
  એવા દરેક પૂર્ણ $\Sigma$-સુસંગત ગણ $\Delta$માં $!A \in \Delta$
  હોય ($\Gamma = \emptyset$ હોય ત્યારે પણ; એ કિસ્સામાં મોડલ તંત્ર
  $\Sigma$નું બીજું લક્ષણચિત્રણ મળે છે).
\end{cor}

\begin{proof}
  ધારો કે $\Gamma \Proves[\Sigma] !A$, અને $\Gamma$ને વિસ્તારે
  એવો કોઈપણ પૂર્ણ $\Sigma$-સુસંગત ગણ $\Delta$ લો. જો $!A
  \notin \Delta$ હોય, તો મહત્તમતાથી $\lnot!A \in \Delta$;
  તેથી એકદિશવર્ધિતાથી $\Delta \Proves[\Sigma] !A$ અને સ્વવાચકતાથી
  $\Delta \Proves[\Sigma] \lnot!A$, એટલે $\Delta$ અસુસંગત છે.
  ઊલટું, જો $\Gamma \Proves/[\Sigma] !A$ હોય, તો $\Gamma \cup \{
  \lnot!A\}$ એ $\Sigma$-સુસંગત છે; લિન્ડનબાઉમના સહાયક પ્રમેયથી
  $\Gamma \cup \{ \lnot!A \}$ને વિસ્તારે એવો પૂર્ણ સુસંગત ગણ
  $\Delta$ છે. સુસંગતતાથી $!A \notin \Delta$.
\end{proof}

\end{document}
